\begin{figure}[ht]
\centering
\begin{tikzpicture}[x=0.42cm, y=0.42cm, font=\footnotesize]

% x-axis: log10 h  from 0 down to -18
% y-axis: log10 |relative error| from 0 down to -18
\draw[->, thick] (-18.5,0) -- (1,0) node[right] {$\log_{10} h$};
\draw[->, thick] (-18.5,-18) -- (-18.5,1) node[above, align=center]
  {$\log_{10}\abs{\text{rel.\ error}}$};

\foreach \k in {0,-3,-6,-9,-12,-15,-18} {
  \draw (\k,0) -- (\k,0.3);
  \node[anchor=south, font=\scriptsize] at (\k,0.35) {$10^{\k}$};
}
\foreach \d in {0,-4,-8,-12,-16} {
  \draw (-18.5,\d) -- (-18.8,\d);
  \node[anchor=east, font=\scriptsize] at (-18.8,\d) {$10^{\d}$};
}

% Centred difference: truncation slope +2 until floor near h ~ 10^{-5.3},
% error ~ 10^{-10.6}, then round-off slope -1.
\draw[thick, engblue] (0,-2) -- (-5.3,-12.6) -- (-15,-3);
\node[engblue, anchor=west, font=\scriptsize] at (-8,-9.2) {centred diff.};

% Richardson-extrapolated centred difference: slope +4 to a lower floor
% h* ~ u^{1/5}, error ~ 10^{-13} then round-off up.
\draw[thick, engamber] (0,-3) -- (-3.5,-16) -- (-9,-9);
\node[engamber!80!black, anchor=west, font=\scriptsize] at (-4.8,-15) {Richardson $O(h^{4})$};

% Complex-step: slope +2 down to full precision then flat at floor (no
% subtraction, so no round-off blow-up); flat at ~10^{-16} for all small h.
\draw[thick, engred] (0,-2) -- (-7,-16) -- (-18,-16);
\node[engred, anchor=west, font=\scriptsize] at (-14,-15.1) {complex step (no floor)};

\draw[thick, dashed, enggray] (-18.5,-16) -- (1,-16);
\node[enggray, anchor=west, font=\scriptsize] at (-3,-17) {binary64 floor};

\end{tikzpicture}
\caption[Finite difference, Richardson, and complex step]{Relative error
of three derivative estimators as the step $h$ shrinks, on log-log axes.
The centred difference falls at slope $+2$ until subtractive round-off
turns it up near $h \approx u^{1/3}$, leaving an attainable error near
$10^{-10.6}$. Richardson extrapolation cancels the leading $h^{2}$ term,
reaching slope $+4$ and a lower floor before its own cancellation wall.
The complex-step estimator $f'(x) \approx \Im[f(x + \mathrm{i} h)]/h$
carries no subtraction of nearly equal reals, so it falls at slope $+2$
straight to the binary64 floor and stays there for any smaller $h$, down
to $h \approx 10^{-200}$. The estimators rank by how they trade
truncation against cancellation, the trade of
Figure~\ref{fig:vol02:ch16:error-vs-h} seen across three algorithms.
Source: authors' analysis; the complex-step method follows
\cite[\S 5.7]{text:higham-numerical}.}
\label{fig:vol02ch16:complex-step}
\end{figure}
